%!TEX root = ../main.tex

\section{Related Work}
\label{sec:related}

\slimparagraph{Knowledge graphs for document-centric RAG.}
Retrieval-augmented generation grounds LLM outputs in external evidence \cite{lewis_retrieval-augmented_2020,guu_realm_2020}. Recent work injects structure by building graphs to aggregate evidence. GraphRAG performs local-to-global retrieval over entity-centric graphs \cite{edge_local_2025}, and related methods retrieve textual subgraphs for multi-hop queries \cite{zhu_knowledge_2025,hu_grag_2025,he_g-retriever_2024,li_simple_2025}. These approaches improve how evidence is organized but the retrieval process itself remains static.

\slimparagraph{Retrieval over semi-structured knowledge bases.}
Complementary to document-centric graph indices, semi-structured knowledge base retrievers directly combine text and explicit relations. KAR grounds query expansion in KG structure \cite{xia_knowledge-aware_2025}, GraphSearch mixes graph and text channels with iterative refinement \cite{yang_graphsearch_2025} and HybGRAG routes queries across a bank of retrieval modules with self-reflection \cite{lee_hybgrag_2025}; CoRAG highlights cooperative hybrid retrieval that preserves global semantic access beyond local neighborhoods \cite{zheng_corag_2025}. In parallel, parametric retrievers such as MoR and mFAR learn to fuse lexical, semantic, and structural signals for ranking \cite{lei_mixture_2025,li_multi-field_2025}. Across these variants, retrieval is framed as scoring candidates from a static index. \retrieverShort differs in that retrieval is formulated as an interactive process: the model dynamically switches between global search and neighborhood expansion, guided by the query requirements and without relying on task-specific supervision.

\slimparagraph{Agents for multi-hop KG retrieval.}
A separate line of work treats the KG as an environment for iterative interaction. Earlier agents follow relation paths using reinforcement learning or learned policies \cite{das_go_2018, xiong_deeppath_2017, sun_pullnet_2019, asai_learning_2020}. In the LLM era, tool-use frameworks such as ReAct \cite{yao_react_2023} and prompt-optimization methods such as AvaTaR \cite{wu_avatar_2024} enable interactive evidence gathering.
Within KG retrieval, traversal-based approaches expand from seed entities using prompted heuristics or learned policies, including Tree-of-Traversals, Think-on-Graph, and GraphFlow \cite{markowitz_tree--traversals_2024,sun_think--graph_2024,yu_can_2025, ma2025think}; related KG-grounded reasoning methods also emphasize multi-step planning or navigation on the KG \cite{luo_reasoning_2024,sun_search--graph_2025}. Traversal-based agents are effective when the correct starting entities are known, but they are prone to anchoring errors and can over-commit to local neighborhoods once exploration begins. In \retrieverShort, global search remains available throughout the trajectory, allowing the agent to retain a complete view of the KG at each step. This design enables coordination between global discovery and multi-hop expansion within a single retrieval. Recent work also explores using reinforcement learning to enable LLMs to explore reasoning paths over KGs~\cite{yan2026explore,wang2026kg}.

Existing work varies in whether it treats retrieval as static ranking over an index or as sequential decision-making on the graph, and in whether it requires graph-specific supervision to learn a ranking function or a traversal policy. \retrieverShort adapts its search strategy online through a minimal, graph-native tool interface. It is training-free; when needed, its tool-use policy can be distilled into compact models from interaction trajectories without ground-truth relevance labels, improving efficiency while preserving retrieval quality.
